%\chapter{Resumo}
\addcontentsline{toc}{chapter}{Resumo}
\label{chap:resumo}


\begin{resumo}

Este projeto propõe o desenvolvimento de algoritmos inovadores baseados em sistemas fuzzy evolutivos para detecção e prevenção de fraudes em sistemas de gerenciamento agropecuário. O diferencial dos sistemas fuzzy evolutivos está na capacidade de adaptar, de forma autônoma, tanto sua estrutura quanto seus parâmetros, à medida que novos dados são recebidos em fluxo contínuo. A modelagem fuzzy permite lidar com incertezas e ambiguidades presentes nos dados reais, utilizando funções de pertinência para representar graus de pertencimento a diferentes padrões ou comportamentos, o que é fundamental para identificar fraudes em cenários complexos e dinâmicos. A crescente ocorrência de fraudes fiscais nesse segmento, evidenciada por operações como CIRA, Rei do Gado e Boi na Linha, reforça a necessidade de abordagens tecnológicas inovadoras. O trabalho será desenvolvido em três etapas: (1) revisão da literatura sobre sistemas inteligentes adaptativos, detecção de anomalias e aprendizado incremental; (2) proposição, desenvolvimento e implementação de algoritmos fuzzy evolutivos, com foco na construção e atualização dinâmica das regras e funções de pertinência; (3) avaliação e validação dos algoritmos propostos em comparação com métodos do estado da arte, utilizando dados reais do setor agropecuário. Espera-se como resultado a obtenção de soluções com alta adaptabilidade, acurácia e interpretabilidade, além de contribuições teóricas e práticas para a detecção de fraudes, com potencial de aplicação em outros domínios de gestão pública.
 
 \vspace{0.5cm}
 
\textbf{Palavras-chave}: Detecção de fraudes em Sistemas, Sistemas Fuzzy Evolutivos, Pertinência, Setor agropecuário.

\end{resumo}
